\documentclass[11pt]{article}
\usepackage[margin=1in]{geometry}
\usepackage{amsmath,amssymb,amsthm}
\usepackage{booktabs}
\usepackage{hyperref}

\title{Full QCD Running of the $b$-Quark Mass\\from $M_P$ to Low Energies\\[6pt]
\large Explicit 2-Loop Computation with Threshold Matching}
\author{DFD Research Programme}
\date{23 March 2026}

\begin{document}
\maketitle

%--------------------------------------------------------------
\section{Setup and Conventions}
%--------------------------------------------------------------

We compute the running of the $b$-quark mass from the Planck scale $M_P = 1.221 \times 10^{19}$~GeV down to $\mu = m_b = 4.18$~GeV (and further to $m_c = 1.27$~GeV), using:
\begin{itemize}
\item 2-loop running of $\alpha_s$ (numerical RK4 integration),
\item 1-loop quark mass anomalous dimension ($\gamma_0 = 8$),
\item Continuous (1-loop) threshold matching at $m_t$, $m_b$, $m_c$.
\end{itemize}

\paragraph{Input parameters (PDG 2024):}
\begin{align}
\alpha_s(M_Z) &= 0.1179, \qquad M_Z = 91.1876~\text{GeV}, \nonumber\\
m_t &= 172.76~\text{GeV}, \quad m_b = 4.18~\text{GeV}, \quad m_c = 1.27~\text{GeV}, \nonumber\\
\alpha_{\mathrm{em}} &= 1/137.036, \qquad v = 246.22~\text{GeV}.
\end{align}

\paragraph{Beta function coefficients.}
For $N_c = 3$ colors and $N_f$ active flavors:
\begin{align}
\beta_0 &= \frac{11 N_c - 2 N_f}{3}, \qquad
\beta_1 = \frac{34 N_c^2 - N_f(10 N_c + 6 C_F)}{3}, \qquad C_F = \tfrac{4}{3}.
\end{align}
Explicitly:
\begin{center}
\begin{tabular}{cccc}
\toprule
$N_f$ & $\beta_0$ & $\beta_1$ & $\gamma_0/(2\beta_0)$ \\
\midrule
6 & 7 & $\tfrac{64}{3} \approx 21.33$ & $4/7 \approx 0.5714$ \\[3pt]
5 & $\tfrac{23}{3} \approx 7.667$ & $\tfrac{116}{3} \approx 38.67$ & $\tfrac{12}{23} \approx 0.5217$ \\[3pt]
4 & $\tfrac{25}{3} \approx 8.333$ & $\tfrac{154}{3} \approx 51.33$ & $\tfrac{12}{25} = 0.4800$ \\[3pt]
3 & 9 & 64 & $4/9 \approx 0.4444$ \\
\bottomrule
\end{tabular}
\end{center}

\paragraph{RG equations.}
The 2-loop running of $\alpha_s$:
\begin{equation}
\mu \frac{d\alpha_s}{d\mu} = -\frac{\beta_0}{2\pi}\,\alpha_s^2 - \frac{\beta_1}{8\pi^2}\,\alpha_s^3.
\end{equation}
The 1-loop mass running:
\begin{equation}
\frac{m(\mu_f)}{m(\mu_i)} = \left[\frac{\alpha_s(\mu_f)}{\alpha_s(\mu_i)}\right]^{\gamma_0/(2\beta_0)},
\qquad \gamma_0 = 8.
\end{equation}

%--------------------------------------------------------------
\section{Results: $\alpha_s$ at All Scales}
%--------------------------------------------------------------

Running $\alpha_s(M_Z) = 0.1179$ upward and downward with 2-loop RGE:

\begin{center}
\begin{tabular}{lcc}
\toprule
Scale & $\alpha_s$ & $N_f$ \\
\midrule
$M_P = 1.22 \times 10^{19}$~GeV & $0.01886$ & 6 \\
$m_t = 172.76$~GeV & $0.1076$ & $6 \to 5$ \\
$M_Z = 91.19$~GeV & $0.1179$ & 5 \\
$m_b = 4.18$~GeV & $0.2232$ & $5 \to 4$ \\
$m_c = 1.27$~GeV & $0.3732$ & $4 \to 3$ \\
\bottomrule
\end{tabular}
\end{center}

Key result: \textbf{$\alpha_s(M_P) \approx 0.0189 \approx 1/53$}, confirming asymptotic freedom drives the coupling to a very small value at the Planck scale.

%--------------------------------------------------------------
\section{Mass Running Factors}
%--------------------------------------------------------------

The running mass ratio $R_i = m(\mu_{\mathrm{low}})/m(\mu_{\mathrm{high}})$ in each region:

\begin{center}
\begin{tabular}{lccc}
\toprule
Region & $N_f$ & Exponent $\gamma_0/(2\beta_0)$ & $R_i$ \\
\midrule
$M_P \to m_t$ & 6 & $4/7 = 0.5714$ & $2.7046$ \\
$m_t \to m_b$ & 5 & $12/23 = 0.5217$ & $1.4636$ \\
$m_b \to m_c$ & 4 & $12/25 = 0.4800$ & $1.2798$ \\
\bottomrule
\end{tabular}
\end{center}

\paragraph{Total running factors:}
\begin{align}
R(M_P \to m_b) &= R_1 \times R_2 = 2.7046 \times 1.4636 = \boxed{3.9583}, \\
R(M_P \to m_c) &= R_1 \times R_2 \times R_3 = 2.7046 \times 1.4636 \times 1.2798 = \boxed{5.0658}.
\end{align}

That is, the $b$-quark mass \emph{increases} by a factor of $\sim 3.96$ when running from $M_P$ down to $m_b$, and by a factor of $\sim 5.07$ when running all the way down to $m_c$.

%--------------------------------------------------------------
\section{Extraction of the Bare Yukawa: $A_b^{\mathrm{bare}}$}
%--------------------------------------------------------------

In the DFD framework, the bare $b$-quark mass at $M_P$ is posited to take the form:
\begin{equation}
m_b(M_P) = A_b^{\mathrm{bare}} \cdot \alpha_{\mathrm{em}} \cdot \frac{v}{\sqrt{2}}.
\end{equation}
The physical (MS-bar) mass at $\mu = m_b$ is then:
\begin{equation}
m_b(m_b) = R(M_P \to m_b) \cdot A_b^{\mathrm{bare}} \cdot \alpha_{\mathrm{em}} \cdot \frac{v}{\sqrt{2}}.
\end{equation}
Solving for $A_b^{\mathrm{bare}}$ with $m_b(m_b) = 4180$~MeV:
\begin{align}
\alpha_{\mathrm{em}} \cdot \frac{v}{\sqrt{2}} &= \frac{1}{137.036} \times 174.104~\text{GeV} = 1.2705~\text{GeV}, \\[6pt]
m_b(M_P) &= \frac{4.180}{3.9583}~\text{GeV} = 1.0560~\text{GeV}, \\[6pt]
A_b^{\mathrm{bare}} &= \frac{m_b(m_b)}{R \cdot \alpha \cdot v/\sqrt{2}}
= \frac{4.180}{3.9583 \times 1.2705} = \boxed{0.8312}.
\end{align}

%--------------------------------------------------------------
\section{Interpretation of $A_b^{\mathrm{bare}} \approx 0.831$}
%--------------------------------------------------------------

The numerical value $A_b^{\mathrm{bare}} = 0.8312$ is strikingly close to several recognizable quantities:

\begin{enumerate}
\item \textbf{$5/6 = 0.8333$} --- agreement to 0.26\%. This is the simplest rational approximation.

\item $A_b^{\mathrm{bare}} \times 12 \approx 9.974$, tantalizingly close to 10.

\item $A_b^{\mathrm{bare}} \times 120 \approx 99.74$, close to 100 (i.e.\ $A_b^{\mathrm{bare}} \approx 100/120 = 5/6$).
\end{enumerate}

If $A_b^{\mathrm{bare}} = 5/6$ exactly, the predicted $m_b(m_b)$ would be:
\begin{equation}
m_b^{\mathrm{pred}}(m_b) = \frac{5}{6} \times 3.9583 \times 1.2705 = 4.190~\text{GeV},
\end{equation}
compared with the PDG value $m_b(m_b) = 4.18 \pm 0.03$~GeV---well within the $1\sigma$ uncertainty.

\paragraph{Yukawa coupling.}
The corresponding Yukawa coupling at $M_P$:
\begin{equation}
y_b(M_P) = A_b^{\mathrm{bare}} \times \alpha_{\mathrm{em}} = 0.8312 \times 0.007297 = 0.006065.
\end{equation}

%--------------------------------------------------------------
\section{Cross-Check: $\alpha_s(M_P)$}
%--------------------------------------------------------------

The 2-loop value $\alpha_s(M_P) = 0.01886$ is consistent with standard results in the literature. For comparison, the 1-loop analytic formula gives $\alpha_s(M_P) = 0.01904$, a 1\% difference that quantifies the 2-loop correction at these extreme scales.

The dominant contribution to the mass running comes from Region~1 ($M_P \to m_t$, $N_f = 6$), where $\alpha_s$ changes by a factor of $\sim 5.7$ over $\sim 39$ decades of energy. This produces a mass enhancement of $R_1 = 2.70$.

%--------------------------------------------------------------
\section{Summary}
%--------------------------------------------------------------

\begin{center}
\renewcommand{\arraystretch}{1.3}
\begin{tabular}{ll}
\toprule
Quantity & Value \\
\midrule
$\alpha_s(M_P)$ & $0.01886 \approx 1/53$ \\
$R(M_P \to m_b)$ & $3.9583$ \\
$R(M_P \to m_c)$ & $5.0658$ \\
$m_b(M_P)$ & $1056$~MeV \\
$A_b^{\mathrm{bare}}$ & $0.8312$ \\
Nearest rational & $5/6 = 0.8333$ ~~(0.26\% deviation) \\
\bottomrule
\end{tabular}
\end{center}

The key finding is that QCD running from $M_P$ to $m_b$ enhances the quark mass by a factor $R \approx 3.96$. In the DFD mass formula $m_b(M_P) = A_b^{\mathrm{bare}} \cdot \alpha \cdot v/\sqrt{2}$, the bare prefactor extracts as $A_b^{\mathrm{bare}} \approx 0.831 \approx 5/6$, a simple rational number that could encode group-theoretic content.

The companion Python script \texttt{full\_QCD\_running\_MP\_to\_1GeV.py} reproduces all numbers in this document and can be modified for sensitivity studies.

\end{document}
